% ══════════════════════════════════════════════════════════════════════════════
%  Chapter 18: Wallet UI (Iced)
% ══════════════════════════════════════════════════════════════════════════════
\chapter{Wallet UI (Iced)}
\label{ch:wallet-iced}

\begin{learnbox}
\begin{itemize}
  \item Build a user-facing wallet interface for managing addresses, balances, and transactions
  \item Implement persistent encrypted storage of wallet data using an embedded database
  \item Extend the Iced MVU pattern with database initialization and query semantics
  \item Design for real wallet workflows: address generation, sending funds, transaction history
\end{itemize}
\end{learnbox}

\lettrine{T}{his} chapter explains the \code{\index{bitcoin-wallet-ui}bitcoin-wallet-ui} \index{Cargo!crate}crate---a \index{user-facing}user-facing \index{wallet}wallet \index{interface}interface for managing \index{address}addresses, \index{balance}balances, and \index{transaction}transactions. Unlike the \index{admin dashboard}admin UI (Chapter~\ref{ch:desktop-iced}), the \index{wallet}wallet UI prioritizes end-user \index{workflow}workflows: creating \index{wallet}wallets, sending \index{transaction}transactions, viewing \index{transaction}transaction \index{history}history, and persisting \index{encryption}encryption \index{key}keys securely.

\begin{notebox}[Focus]
This chapter focuses on the \index{wallet}wallet \index{interface}interface, not the \index{wallet}wallet \index{implementation}implementation itself. The \index{wallet}wallet \index{backend}backend is explained in Chapter~\ref{ch:wallet} (\index{Wallet}Wallet \index{System}System). Here we build the \index{MVU}MVU shell and \index{persistence}persistence layer that puts that \index{wallet}wallet to work in a user's hands.
\end{notebox}

The \index{wallet}wallet UI uses the same \index{Iced}Iced \index{MVU}MVU \index{pattern}pattern as Chapter~\ref{ch:desktop-iced} but adds \textbf{\index{persistent}persistent \index{encryption}encrypted \index{storage}storage}. An \index{embedded database}embedded \index{database}database (Chapter~\ref{ch:embedded-db}) stores \index{wallet}wallet \index{metadata}metadata and \index{derived state}derived \index{data}data, eliminating the need to re-query the \index{blockchain}blockchain for every \index{session}session.

\section{What This UI Is (In One Sentence)}
\label{sec:ch18-oneline}

The wallet UI is a \textbf{user-facing MVU application} that adds \textbf{persistent encrypted storage}: it renders wallet state (\code{view.rs}), mutates state in response to user actions (\code{update.rs}), performs async API calls (\code{api.rs}), and persists balances and transaction history to an encrypted database (\code{database.rs}).

\section{Architectural Map}
\label{sec:ch18-arch-map}

This wallet UI is a message-driven state machine that adds persistence to the MVU pattern:

\begin{minted}[fontsize=\footnotesize]{text}
View
  view(&WalletApp) -> Element<Message>
  |
  +-- User action produces Message
  |
  v
Update
  update(&mut WalletApp, Message) -> Task<Message>
  |
  +-- Mutates model (WalletApp)
  +-- Returns Task for async work
  |
  v
Model
  WalletApp
  +-- in-memory state (addresses, balances, tx history)
  +-- references to database for persistence
  |
  v
Task<Message>
  +-- spawn_on_tokio(...) for async work
  |
  v
Tokio bridge
  +-- Database queries (Chapter 29)
  +-- API calls to WalletClient HTTP
  |
  v
Persistent Storage + Node API
  +-- Encrypted database (SQLite with XChaCha20)
  +-- Bitcoin Node API (Chapter 23)
\end{minted}

% ──────────────────────────────────────────────────────────────────────────────
\section{Boot Sequence and Wiring}
\label{sec:ch18-boot}

\begin{listing}[H]
\caption{Main entrypoint (\code{bitcoin-wallet-ui/src/main.rs})}
\label{lst:ch18-main}
\begin{minted}[fontsize=\footnotesize]{rust}
mod api;
mod app;
mod database;
mod runtime;
mod types;
mod update;
mod view;

use app::WalletApp;
use iced::{Theme, application};
use runtime::init_runtime;
use update::update;
use view::view;

fn main() -> iced::Result {
    // 1) Async precondition: initialize Tokio runtime for API calls
    init_runtime();

    // 2) Persistence precondition: generate deterministic database password
    let db_password = generate_database_password();
    if let Err(e) = database::init_database(&db_password) {
        eprintln!("Failed to initialize database: {}", e);
    }

    // 3) Enter Iced's MVU loop
    application(WalletApp::new, update, view)
        .title(title)
        .theme(theme)
        .run()
}

fn generate_database_password() -> String {
    use std::collections::hash_map::DefaultHasher;
    use std::hash::{Hash, Hasher};

    let mut hasher = DefaultHasher::new();

    // Deterministic password from username + home directory
    if let Ok(username) = std::env::var("USER") {
        username.hash(&mut hasher);
    } else if let Ok(username) = std::env::var("USERNAME") {
        username.hash(&mut hasher);
    }

    if let Some(home) = dirs::home_dir() {
        home.to_string_lossy().hash(&mut hasher);
    }

    // Return hex-encoded hash as password
    format!("{:x}", hasher.finish())
}
\end{minted}
\end{listing}

\textbf{Walkthrough: three preconditions before MVU starts:}

\begin{enumerate}
  \item \textbf{Async execution}: \code{init\_runtime()} (same as Chapter~\ref{ch:desktop-iced}) creates a Tokio runtime for API calls.
  \item \textbf{Persistent storage}: \code{init\_database(password)} opens/creates an encrypted SQLite database (see Chapter~\ref{ch:embedded-db}).
  \item \textbf{MVU loop}: \code{application(WalletApp::new, update, view).run()} starts the event loop.
\end{enumerate}

The database password is generated deterministically from the user's username and home directory, so:
\begin{itemize}
  \item Same user on same machine: same password, so the database can be reopened.
  \item Different user or machine: different password, so wallets are not trivially portable.
\end{itemize}

% ──────────────────────────────────────────────────────────────────────────────
\section{Wallet Application State}
\label{sec:ch18-wallet-state}

\begin{listing}[H]
\caption{Wallet state (\code{bitcoin-wallet-ui/src/app.rs})}
\label{lst:ch18-app}
\begin{minted}[fontsize=\footnotesize]{rust}
use bitcoin_api::WalletClient;

#[derive(Debug)]
pub struct WalletApp {
    // Current wallet
    pub current_wallet: Option<String>,
    pub wallets: Vec<String>,

    // Balances and transaction history
    pub balance: Option<u64>,
    pub transactions: Vec<Transaction>,
    pub unspent_outputs: Vec<UnspentOutput>,

    // UI state
    pub active_tab: Tab,
    pub selected_address: Option<String>,
    pub input_send_to: String,
    pub input_amount: String,
    pub status: String,

    // Derivation path management
    pub derivation_path: String,
    pub generated_addresses: Vec<String>,
}

#[derive(Debug, Clone)]
pub struct Transaction {
    pub txid: String,
    pub timestamp: i64,
    pub amount: i64,
    pub fee: Option<i64>,
    pub from_addresses: Vec<String>,
    pub to_addresses: Vec<String>,
}

#[derive(Debug, Clone)]
pub struct UnspentOutput {
    pub txid: String,
    pub vout: usize,
    pub amount: u64,
    pub address: String,
}

#[derive(Debug, Clone, Copy, PartialEq, Eq)]
pub enum Tab {
    Addresses,
    SendTransaction,
    History,
    Balance,
}

impl Default for WalletApp {
    fn default() -> Self {
        Self {
            current_wallet: None,
            wallets: Vec::new(),
            balance: None,
            transactions: Vec::new(),
            unspent_outputs: Vec::new(),
            active_tab: Tab::Addresses,
            selected_address: None,
            input_send_to: String::new(),
            input_amount: String::new(),
            status: String::from("Ready"),
            derivation_path: "m/44'/0'/0'/0".to_string(),
            generated_addresses: Vec::new(),
        }
    }
}

impl WalletApp {
    pub fn new() -> (Self, iced::Task<Message>) {
        (Self::default(), iced::Task::none())
    }
}
\end{minted}
\end{listing}

\textbf{Walkthrough: wallet state design:}

\code{WalletApp} groups state into three categories:

\begin{itemize}
  \item \textbf{Wallet identity}: \code{current\_wallet} (active wallet address), \code{wallets} (list of known wallets).
  \item \textbf{Derived data}: \code{balance}, \code{transactions}, \code{unspent\_outputs} (cached from API or database).
  \item \textbf{UI state}: \code{active\_tab}, \code{input\_*} (user inputs), \code{status} (feedback message).
\end{itemize}

Unlike the admin UI, the wallet UI is designed around a \textbf{single active wallet}---the user selects one wallet to manage at a time.

% ──────────────────────────────────────────────────────────────────────────────
\section{Message Types}
\label{sec:ch18-messages}

\begin{listing}[H]
\caption{Message enum (excerpt) (\code{bitcoin-wallet-ui/src/types.rs})}
\label{lst:ch18-messages}
\begin{minted}[fontsize=\footnotesize]{rust}
#[derive(Debug, Clone)]
pub enum Message {
    // Wallet selection
    SelectWallet(String),
    FetchWalletList,
    WalletListLoaded(Result<Vec<String>, String>),

    // Balance and transaction history
    FetchBalance(String),
    BalanceLoaded(Result<BalanceResponse, String>),

    FetchTransactionHistory(String),
    HistoryLoaded(Result<Vec<Transaction>, String>),

    FetchUnspentOutputs(String),
    UnspentLoaded(Result<Vec<UnspentOutput>, String>),

    // Address generation
    GenerateAddress,
    AddressGenerated(Result<String, String>),

    // Sending transactions
    SendToChanged(String),
    AmountChanged(String),
    SendTransaction,
    TransactionSent(Result<String, String>),

    // Navigation
    TabChanged(Tab),

    // Persistence
    SaveToDatabaseRequested,
    SaveToDatabase(Result<(), String>),
}
\end{minted}
\end{listing}

Messages fall into four categories:

\begin{itemize}
  \item \textbf{Wallet operations}: select wallet, fetch wallets, generate addresses.
  \item \textbf{Balance/history queries}: fetch balance, transaction history, unspent outputs.
  \item \textbf{Transactions}: send Bitcoin (user inputs and submission).
  \item \textbf{Persistence}: save and load from encrypted database.
\end{itemize}

% ──────────────────────────────────────────────────────────────────────────────
\section{Update Loop}
\label{sec:ch18-update}

\begin{listing}[H]
\caption{Update loop (excerpt) (\code{bitcoin-wallet-ui/src/update.rs})}
\label{lst:ch18-update}
\begin{minted}[fontsize=\footnotesize]{rust}
use crate::api;
use crate::app::WalletApp;
use crate::database;
use crate::runtime::spawn_on_tokio;
use crate::types::Message;
use iced::Task;

pub fn update(app: &mut WalletApp, message: Message) -> Task<Message> {
    match message {
        Message::TabChanged(tab) => {
            app.active_tab = tab;
            Task::none()
        }

        Message::SelectWallet(wallet) => {
            app.current_wallet = Some(wallet.clone());
            app.status = format!("Selected wallet: {}", wallet);
            // Automatically fetch balance for the selected wallet
            let task_fetch_balance = Task::perform(
                spawn_on_tokio(api::fetch_balance(wallet.clone())),
                Message::BalanceLoaded,
            );
            task_fetch_balance
        }

        Message::FetchBalance(wallet_addr) => {
            Task::perform(
                spawn_on_tokio(api::fetch_balance(wallet_addr)),
                Message::BalanceLoaded,
            )
        }

        Message::BalanceLoaded(result) => {
            match result {
                Ok(response) => {
                    app.balance = Some(response.balance);
                    app.status = format!("Balance: {} satoshis", response.balance);
                    // Persist balance to database
                    if let Some(ref wallet) = app.current_wallet {
                        let wallet = wallet.clone();
                        let balance = response.balance;
                        Task::perform(
                            spawn_on_tokio(async move {
                                database::save_balance(&wallet, balance).await
                            }),
                            |result| Message::SaveToDatabase(result),
                        )
                    } else {
                        Task::none()
                    }
                }
                Err(e) => {
                    app.status = format!("Error fetching balance: {}", e);
                    Task::none()
                }
            }
        }

        Message::SendToChanged(address) => {
            app.input_send_to = address;
            Task::none()
        }

        Message::AmountChanged(amount) => {
            app.input_amount = amount;
            Task::none()
        }

        Message::SendTransaction => {
            if let Some(ref current_wallet) = app.current_wallet {
                if app.input_send_to.is_empty() || app.input_amount.is_empty() {
                    app.status = "Recipient and amount are required".to_string();
                    return Task::none();
                }

                let from = current_wallet.clone();
                let to = app.input_send_to.clone();
                let amount: u64 = match app.input_amount.parse() {
                    Ok(a) => a,
                    Err(_) => {
                        app.status = "Invalid amount".to_string();
                        return Task::none();
                    }
                };

                Task::perform(
                    spawn_on_tokio(api::send_transaction(from, to.clone(), amount)),
                    Message::TransactionSent,
                )
            } else {
                app.status = "No wallet selected".to_string();
                Task::none()
            }
        }

        Message::TransactionSent(result) => {
            match result {
                Ok(txid) => {
                    app.status = format!("Transaction sent: {}", txid);
                    app.input_send_to.clear();
                    app.input_amount.clear();
                }
                Err(e) => {
                    app.status = format!("Error sending transaction: {}", e);
                }
            }
            Task::none()
        }

        _ => Task::none(),
    }
}
\end{minted}
\end{listing}
